\section{Introduction}

During the last 20 years chess has become a useful empirical environment for studying high-skill decision making under time constraints. Chess is attracting increasing attention from data scientists and economists as a great source of data. Unlike many real-world production settings, chess also allows the researcher to evaluate the quality of individual decisions (moves) using strong chess engines. These features make elite chess a rare case in which individual characteristics, incentives, rule design, and decision quality can be observed in the same empirical setting. In addition, thanks to changes in the tournament rules, we can see how the rules themselves affect the final results.

This thesis studies a major rule change in Chess.com Titled Tuesday tournament. Titled Tuesday is a weekly online blitz tournament restricted to titled chess players. The event is the largest regular online tournaments for elite players and includes grandmasters, international masters, national masters, and other officially titled by FIDE players. Chess.com describes the modern event as an 11-round Swiss tournament for verified titled players with cash prizes and special prizes for women and streamers \citep{chesscom_titled_tuesday}. The tournament is not a laboratory experiment. It is a recurring high-skill online competition in which players repeatedly decide whether to participate and repeatedly face opponents of comparable strength. To emphasize the importance of this tournament, it is also worth noting that Titled Tuesday became part of the Champions Chess Tour and Esports World Cup cycle \citep{chesscom_titled_tuesday}.

The institutional change studied here took place at the start of September 2025. Before the reform, the post-February-2022 Titled Tuesday format consisted of two weekly tournaments, commonly referred to as early (9:00 a.m. US Eastern Time) and late (5:00 p.m. US Eastern Time) events, with a 3+1 blitz time control: three minutes for the game for each player and one second increment after each move. Chess.com announced that the new Titled Tuesday season would start on September 2, 2025 with a single weekly tournament and a 5+0 time control: five minutes for the game and no increment \citep{chesscom_major_changes_2025}. In other words, we are actually seeing two changes: 1) a change in the time controls in chess, and 2) a change in the schedule - instead of two tournaments for different zones, there is now a single tournament for everyone.  

Both parts of the reform are actually economically important changes. First, it replaced incremental blitz with no-increment blitz. This changes the time-budget technology of the game. In a 3+1 game, a player who survives into a long endgame receives one extra second per move, which lowers the risk of losing entirely because of clock exhaustion. In a 5+0 game, the player starts with a larger initial time budget, but the marginal time cost of every move is final: no time is ever added back. In addition, since the average game consists of far fewer than 120 moves, players now have significantly more time on average following the reform. Second, the reform removed the old second time slot. This changed the participation environment because players who were regular late-slot participants could no longer choose their preferred time window.

\begin{tcolorbox}[
    colback=white,
    colframe=blue,
    boxrule=0.6pt,
    arc=1mm,
    left=8pt,
    right=8pt,
    top=6pt,
    bottom=6pt
]

The core research question is:


\textit{How did the switch from 3+1 to 5+0 affect performance, participation, and adaptability among elite online chess players?}
\end{tcolorbox}

The analysis uses several datasets that have been combined into a single database. The main player-game panel covers Titled Tuesday games from February 2022\footnote{We are limiting our analysis to the period starting from that date, as the previous rule change took place in February 2022. At that time, the format shifted from one tournament per week to two tournaments, the prize pool was increased, with special prizes added for women and streamers.} through April 2026. It contains player names, opponent names, ratings, titles, color of pieces, round, event identifiers, game outcomes (results, who won in a game), accuracy measures, tournament information, and demographic or country-linked variables where available. The move-level part of the thesis uses Chess.com PGN records, per-move clock annotations, and Stockfish engine evaluations. The current Stockfish mechanism sample contains 27,517,197 usable move-level observations, 319,095 games, 616,731 player-games, 5,768 players, and 148 tournaments. The clock-augmented mechanism sample is larger and contains about 1.56 million matched player-games and 69.0 million weighted move observations. In addition, in this work we construct player-level pre-change style features from move data and use a neural contrastive learning model to create player style embeddings from pre-change games.

The empirical design is not a randomized experiment. The identifying variation comes from a sharp institutional break in September 2025 and from repeated observations of the same players before and after the change. The main game-level specifications use player fixed effects and event fixed effects, with controls for rating difference, color, and tournament round. Robustness checks include excluding early rounds, restricting the time window around the reform, using balanced-player samples, paired-game fixed effects, placebo rule-change dates, event-study specifications, and two-way clustered standard errors. Move-level specifications use player fixed effects, tournament fixed effects, and move-number or phase controls. The clock specifications additionally use remaining-time thresholds, time spent on the move, first entry into low-clock states, and clock-evaluation snapshots at moves 10, 20, and 30.

The main game-level result is that the reform increased the return to relative skill. After the switch to 5+0, a 100-point rating advantage became worth about 1.2 additional percentage points of score share. The estimate is stable across event fixed effects, balanced-player samples, near-window samples, exclusion of the first two rounds, and paired-game fixed effects. Placebo cutoff estimates are much smaller. Win/loss decomposition shows that the effect comes from higher win probability and lower loss probability for stronger players, with little change in draw probability. This pattern is inconsistent with the simple claim that no-increment chess mainly increases randomness.

The second main result is that move quality became more tightly linked to outcomes. A 10-point accuracy advantage converted into about 1.4 additional percentage points of score share in the cleaned thesis table, and a broader production-function battery estimates a larger coefficient of about 3.5 percentage points. Since accuracy is measured after the game, this is not interpreted as a causal effect of accuracy. The result instead describes a steeper post-change production function: when one player outplayed the opponent by the same measured quality margin, that advantage translated more strongly into points under 5+0.

The third result concerns heterogeneity in adaptation. Young players performed better after this reform. The main age interaction is about -1.5 percentage points of score share per additional 10 years of age. Within-game age-gap specifications and young-versus-old paired comparisons point in the same direction, although placebo checks show that some age-specific trend existed before the reform. Players with strong Chess.com ratings relative to classical ratings also gained, suggesting a role for online-platform capital: interface fluency, premove skill, mouse speed, and practical experience in the Chess.com environment. This online-capital result is robust as an association but less clean as a sharp discontinuity because pretrend checks indicate a broader trend.

The fourth result separates the time-control effect from the schedule component. Old late-slot regular players became less likely to appear in post-reform events and completed fewer rounds conditional on returning. A pure late-slot player was about 13.3 percentage points less likely to appear after the reform and played about 1.38 fewer post-reform tournaments out of 31 available events. Among late regulars who did return, tournament score and rank fell, but per-game accuracy and per-game result did not clearly decline. The schedule shock therefore appears to operate mainly through availability and completion, not through lower move quality conditional on playing.

The move-level and clock results explain how the skill premium appears inside games. Late-game and final-10-ply centipawn losses increased under 5+0. Equal or critical positions became more costly, with higher capped centipawn loss and higher blunder probability. The clock data show that ordinary low-clock exposure became less frequent after the reform, but severe time trouble became more damaging: the post-change penalty is concentrated below 10 seconds, and recovery from very low clock states almost disappears without increment. Rating favorites converted winning positions more effectively after the reform, converted +2 positions faster, were less likely to let advantages drop back to equality, and recovered more from -2 positions. Clock weakness also changed conversion directly: being personally low on time while winning became more costly, while facing an opponent already below 10 seconds became more valuable for conversion and defensive escape. Pre-change defensive skill and complexity tolerance predicted lower post-change error rates. These findings support a conversion-technology interpretation: 5+0 increased the value of turning advantages into wins and of surviving inferior positions under harsher clock constraints.

Finally, the thesis studies player style. Interpretable pre-change style indexes classify players along clean-play, chaos-creation, self-risk, and practical-chaos dimensions. The econometric style-index results show that historically chaotic profiles did not receive a relative score premium from the new format; if anything, one-standard-deviation higher chaos style is associated with about one percentage point lower post-change score share in several specifications. A neural contrastive style model trained on pre-change move sequences also produces meaningful player clusters. These clusters show broad raw accuracy gains after the reform across clean, average, lower-score, and volatile style groups, which suggests that the move-quality gains from the extra initial time were not confined to one narrow player type.

The contribution of the thesis is threefold. First, it provides evidence on adaptability to a real time-constraint shock in a high-skill competitive environment. Second, it links game-level panel econometrics to move-level computational measurement, allowing the same rule change to be studied through outcomes, participation, centipawn loss, conversion, and style. Third, it adds a data-science layer by constructing interpretable and neural measures of player style from pre-change behavior.

The rest of the thesis is organized as follows. Section 2 reviews related literature on chess as an empirical setting for economics, cognitive performance, time pressure, gender, and machine learning. Section 3 describes the data sources, scraping process, cleaning rules, and constructed variables. Section 4 presents the empirical strategy. Section 5 reports the main robust results. Section 6 discusses analyses that were run but did not produce robust evidence. Section 7 concludes and discusses limitations and future work.
